%% elsarticle LaTeX template
%% Transportation Research Part C: Emerging Technologies
%% Submission namespace: fx2_pop1_a800_t500_dc150_r3
%%
%% All numerical claims are from the locked namespace above (10 paired seeds,
%% N=100 customers, pop-weighted demand, min-of-3 evaluation budget).

\documentclass[review,3p,authoryear]{elsarticle}

\usepackage{amsmath,amssymb,amsthm}
\usepackage{booktabs}
\usepackage{multirow}
\usepackage{array}
\usepackage{graphicx}
\graphicspath{{../figures/}{figures/}}
\usepackage{xcolor}
\usepackage{hyperref}
\hypersetup{hidelinks,hypertexnames=false}
\usepackage{natbib}
\usepackage{microtype}
\usepackage{siunitx}
\usepackage{algorithm}
\usepackage{algpseudocode}
\usepackage[capitalize,noabbrev]{cleveref}
\usepackage{enumitem}

% theorem environments
\newtheorem{theorem}{Theorem}[section]
\newtheorem{lemma}[theorem]{Lemma}
\newtheorem{proposition}[theorem]{Proposition}
\newtheorem{corollary}[theorem]{Corollary}
\newtheorem{definition}[theorem]{Definition}
\newtheorem{remark}[theorem]{Remark}

% notation shorthands
\newcommand{\Rstar}{R^{*}}
\newcommand{\Qprime}{Q^{\prime}}
\newcommand{\eff}{\mathrm{eff}}
\newcommand{\SoC}{\mathrm{SoC}}
\newcommand{\cP}{\mathcal{P}}
\newcommand{\cS}{\mathcal{S}}
\newcommand{\cR}{\mathcal{R}}
\newcommand{\ELRP}{\text{GW-ELRP}}

\journal{Transportation Research Part C: Emerging Technologies}

%%----------------------------------------------------------------------
\begin{document}

\begin{frontmatter}

\title{Battery size moderates the value of terrain and winter information\\
in charger siting for regional electric freight}

%% Authors — fill in before submission
\author[trond]{Surendra Reddy Kancharla}
\ead{kanna.srk@icloud.com}

\address[trond]{[Affiliation -- to be completed]}

%%----------------------------------------------------------------------
\begin{abstract}
We study how omitting road grade or winter energy losses changes
charger-siting decisions for regional electric freight.  The model combines
arc-level grade, payload-dependent traction and regenerative braking,
nonlinear charging, and budget-constrained site selection.  A conservative
state-of-charge-grid dynamic program evaluates fixed routes; adaptive
large-neighbourhood search generates a route pool assembled by a
set-partitioning mixed-integer program.  On 44 small flat-network benchmark
instances, the matheuristic agrees with a discretized state-space reference
on 42 (median gap 0.00\%, mean 0.25\%).  A further 168 instances with up to
320 customers test compatibility but are not quality-certified against
best-known solutions.

A $2\times2$ experiment compares full, terrain-only, winter-only, and blind
planners on ten paired 100-customer Trondheim instances, all with a free
150\,kW depot charger.  Averaged over budgets $B=2,\ldots,8$, blind planning
adds an estimated 17.5\,min/day for a 145\,kWh pack (95\% confidence interval
0.8--34.2) and 60.8\,min/day for a 100\,kWh pack (16.6--104.9), equivalent
to 0.86\% and 2.97\% of baseline route duration.  Component means change
ordering with battery size, but paired terrain--winter contrasts include
zero.  Regret shows no decreasing budget trend.  Blind planners use fewer
than one non-depot site on average even at $B=8$, indicating that omitted
energy terms suppress perceived charging demand.  No customer is left
unserved.

A continuum threshold identifies 100\,kWh winter as Trondheim's only
mean-tour capacity exceedance.  Transfer calculations for Stuttgart and
Groningen expose sensitivity to digital-surface-model bias.  The evidence
supports battery-dependent information value, but not a universal
terrain-versus-winter dominance rule.
\end{abstract}

\begin{keyword}
Electric vehicle routing \sep charger siting \sep road grade \sep
regenerative braking \sep winter range degradation \sep placement regret \sep
matheuristic \sep fixed-route charging problem \sep value of information
\end{keyword}

\end{frontmatter}

%%======================================================================
\section{Introduction}
\label{sec:intro}

Freight decarbonisation is placing battery-electric trucks on road
networks that are neither flat nor seasonally uniform.  Field experience
shows that the operational viability of these fleets hinges on
route-specific operating profiles and battery
sizing~\citep{Link2022,Hovi2019}.  Planners who site en-route chargers
must therefore decide how much to invest in characterising the road
network's topography and the seasonal variation in vehicle consumption.
Many routing and siting studies represent energy use by a constant
distance rate or a scenario-independent range.  Whether that simplification
changes infrastructure decisions, rather than only route-energy estimates,
is less well established.

One practical objection is that depot charging makes en-route siting moot.
A truck
that can swing back to its depot to charge mid-tour has a natural
alternative to any public en-route charger.  In a hub-and-spoke urban
distribution network the depot charger is effectively free infrastructure,
so en-route charger quality matters only at the margin.

We test that objection directly and separate the information effects with
a factorial design.  The main result is a battery-dependent difference in
the value of model fidelity.  A blind planner has a substantially larger
route-duration penalty with a 100\,kWh pack than with a 145\,kWh pack.
The component means change ordering across the two packs, although ten
paired instances do not estimate that contrast precisely enough to support
a categorical ``dominance flip.''  Regret also shows no decreasing trend as
the site-count upper bound rises from two to eight.  This occurs because the
blind model generally leaves the additional allowance unused, not because
extra sites are proven incapable of correcting a bad ranking.  Every planner
has access to a free 150\,kW depot charger, and no evaluated cell leaves a
customer unserved.

\subsection*{Magnitude framing}

Across budgets $B=2,\ldots,8$, the estimated penalty of a fully blind
planner relative to the full-information heuristic baseline averages
60.8\,min/day for a 100\,kWh fleet and 17.5\,min/day for a 145\,kWh
fleet.  These are sums of route completion times, so they can be read as
driver-minutes only under the assumption that labour cost is proportional
to route duration.

\subsection*{Contributions}

On the modelling and algorithmic side, we formulate the grade+winter
ELRP (\ELRP): an electric location-routing problem with
arc-level grade-dependent consumption, regenerative braking with
battery-saturation clamping, a seasonal energy factor, and a
budget-constrained site-selection layer.  We implement a conservative
SoC-grid FRVCP dynamic program for the graded energy model and show that
the energy representation remains piecewise linear under regeneration and
battery saturation.  The usual load-monotone dominance argument cannot be
used unchanged on regen-dominant arcs (35.9\% of valid Trondheim arcs), so
the small-instance reference method retains the delivered-customer subset
explicitly.  The \ELRP{} is solved with an ALNS route-pool matheuristic and
a set-partitioning MILP.  We run the flat special case on 212 Schneider
(2014) and Montoya (2017) instances with up to 320 customers; 42 of 44
small instances numerically agree with the discretized reference
(0.00\% median and 0.25\% mean gap).

On the analytical side, we derive a closed-form ruggedness threshold
$\Rstar$ that screens which battery $\times$ season combinations have mean
tour energy above usable capacity, and assess prospectively recorded
predictions on Stuttgart and Groningen.  Finally, through a
2$\times$2 factorial regret experiment with common-random-number paired
seeds, we quantify how total and component information effects change with
battery size, test whether estimated regret declines with the charger-budget
upper bound, and separate route-duration penalties from unserved-customer
outcomes under depot charging.

\subsection*{Scope}

This paper provides single-region siting experiments for Trondheim (Norway).
Stuttgart and Groningen are transfer checks for the aggregate threshold, not
full siting replications.  The threshold is a mean-tour screening relation,
not an optimization-derived density or a prediction of the discrete siting
optimum.  Cross-regional causal or policy generalisation is therefore outside
the evidence presented here.

\subsection*{Paper organisation}

\Cref{sec:related} reviews the relevant literature.
\Cref{sec:model} formalises the road network, energy model, and ruggedness
threshold.
\Cref{sec:elrp} defines the \ELRP{} and its algorithmic components.
\Cref{sec:validation} validates the solver on standard benchmarks.
\Cref{sec:experiment} describes the factorial regret experiment.
\Cref{sec:results} presents the results.
\Cref{sec:discussion} discusses implications and limitations.
Appendices provide the PWL monotonicity argument, the threshold derivation,
and details of the conservative SoC discretization.

%%======================================================================
\section{Related Work}
\label{sec:related}

\subsection{Electric vehicle routing and the location-routing problem}

The electric vehicle routing problem (E-VRP) has generated a substantial
body of research since the introduction of the first E-VRPTW
formulation~\citep{Schneider2014}.  Pelletier et al.~\citep{Pelletier2016} provide an
early synthesis of the modelling building blocks for electric goods
distribution, covering battery constraints, charging decisions, and
infrastructure dependencies, and argue that siting and routing cannot be
decoupled once realistic range limitations are imposed.  More recent
surveys~\citep{Erdelic2019,Kucukoglu2021,Qin2021} classify the
rapidly-growing literature into families by charging policy (full,
partial, opportunity), fleet composition, and energy model, and
collectively cover over 160 papers on E-VRP variants.  The problem of jointly
optimising charger locations and vehicle routes---the electric
location-routing problem (ELRP)---is reviewed by
Schiffer et al.~\citep{Schiffer2019rev}, who map its connections to the broader class of
routing problems with intermediate stops and show that even moderate
fleet sizes render the joint problem computationally demanding.
Across these surveys, energy consumption is usually represented by a
constant distance rate or by charging-state dynamics; arc-level terrain and
seasonal energy factors are uncommon in the reviewed location-routing
formulations.

\subsection{Nonlinear charging and the fixed-route charging problem}

The most consequential modelling step in recent E-VRP research is the
transition from linear to nonlinear (concave) charging
functions~\citep{Montoya2017}.  Montoya et al.~\citep{Montoya2017} show that forcing
linear charging on nonlinear hardware can yield routes that are infeasible
or materially suboptimal when evaluated on real charging curves, and
they introduce the fixed-route vehicle charging problem (FRVCP) as the
key subproblem: given a fixed customer sequence, find the optimal charging
schedule.  Froger et al.~\citep{Froger2019} develop strengthened formulations and an
efficient labeling algorithm for the FRVCP based on piecewise-linear
SoC-function labels, which underpins the open-source solver
frvcpy~\citep{Kullman2021joc} and several subsequent exact and heuristic
algorithms~\citep{Lee2021,Kullman2021ts}.  Froger et al.~\citep{Froger2021cap}
extend the E-VRP-NL with capacitated charging stations, pairing an
iterated-local-search route generator with an exact route-assembly
phase---the same generate-then-assemble pattern adopted here.  Pelletier et al.~\citep{Pelletier2018} apply
FRVCP-style DP to charge-scheduling for fixed-route electric freight fleets
over multiple days, demonstrating that charging decisions can materially affect
fleet operating cost when battery limits interact with service patterns.
The cited formulations do not jointly propagate arc-level terrain and a
seasonal energy scenario into infrastructure location.

\subsection{Grade and road topology in EV energy models}

The effect of road grade on EV energy consumption is well-established in
the vehicle-engineering literature.  Longitudinal-dynamics models express
arc energy as the sum of rolling resistance, aerodynamic drag, and a
gravitational term proportional to elevation change and
mass~\citep{Fiori2016}.  Basso et al.~\citep{Basso2019} integrate such a model
directly into the E-VRP and demonstrate that neglecting elevation changes
leads to systematic underestimates of energy use and downstream routing
infeasibilities.  Data-driven approaches using GPS and CAN-bus data
confirm the empirical importance of grade for route-level range
prediction~\citep{DeCauwer2017}, and optimisation studies show that
terrain-aware route selection can reduce battery draw even
before charger placement is considered~\citep{Perger2020}.  These works,
however, primarily treat grade as an input to routing.  Comparatively little
work asks how grade changes the \emph{siting} of charger infrastructure or
quantifies the cost of omitting grade at siting time.

\subsection{Cold-weather range degradation}

At low ambient temperatures, usable battery capacity decreases as
internal resistance rises, regenerative braking efficiency falls, and
auxiliary heating loads grow substantially~\citep{Neubauer2014}.
Steinstr\"{a}ter et al.~\citep{Steinstraeter2021} decompose observed cold-weather
range loss into heating demand (up to 31.9\%) and limited regeneration
(up to 21.7\%), with a combined maximum near 50\% in their passenger-car
data.
Laboratory dynamometer studies~\citep{AAA2019} confirm that at
$-6.7\,^{\circ}$C with HVAC enabled, range falls on average by 41\%.
Yuksel and Michalek~\citep{Yuksel2015} show that cold climate regions incur roughly a 24\%
efficiency penalty attributable primarily to cabin heating.  For
commercial freight in particular, field experience from Norway---the
country of our primary study city---documents winter range reductions and
operational adaptations in electric truck fleets~\citep{Hovi2019}.
We did not identify a prior ELRP study that isolates the placement effect of
omitting a seasonal energy factor while holding the demand instances and
evaluation model fixed.

\subsection{Charging infrastructure location models}

Schiffer and Walther~\citep{Schiffer2017} introduce the ELRP with time windows and partial
recharging as the foundational model for joint charger-siting and routing,
and Schiffer and Walther~\citep{Schiffer2018} extend it to robust strategic network design under
operational uncertainty, using an ALNS to generate candidate routing
structures.  Kchaou-Boujelben~\citep{Kchaou2021} surveys the broader charging-station
location literature and documents the diversity of objectives, demand
representations, and range assumptions in that field.  He et al.~\citep{He2018} demonstrate
that even in a network-flow siting model, the assumed driving range
changes optimal station placement---a result that motivates
asking how \emph{energy-model misspecification} (rather than range
uncertainty) propagates to siting decisions.

\subsection{Value of information and placement regret}

Uncertainty in facility location has been studied primarily from the
demand side: stochastic-programming formulations measure the expected cost
of siting under unknown customer patterns, and regret-based models hedge
against scenario-worst outcomes~\citep{Snyder2006}.  Within the E-VRP
literature, Pelletier et al.~\citep{Pelletier2019} quantify the operational consequences of
energy-consumption uncertainty on routing, finding that simplified energy
assumptions can produce routes that are infeasible or expensive when
evaluated under a richer model.  Our work is complementary: rather than
stochastic demand or uncertain consumption, we study \emph{deterministic
model misspecification}---a planner who simply omits terrain or season from
the energy model---and measure the resulting placement regret under the
full-information evaluation model using a controlled factorial design with
common-random-number variance reduction~\citep{Glasserman1992}.
Our specific contribution is to estimate separate terrain and winter
information effects within one charger-siting pipeline.

\subsection{Metaheuristics, matheuristics, and continuum approximation}

The ALNS framework of Ropke and Pisinger~\citep{Ropke2006} is the standard methodological
anchor for rich vehicle routing; Schiffer et al.~\citep{Schiffer2019alns} apply an ALNS
to the location-routing problem with intraroute facilities, the closest
prior algorithmic precedent for our charger-siting matheuristic.  The
pattern of generating route pools via local search and assembling them
with a set-partitioning MILP is a recognised family of VRP
matheuristics~\citep{Archetti2014,Montoya2017}.  For the strategic
threshold result, the continuum-approximation framework of
Daganzo~\citep{Daganzo1984} provides closed-form tour-length estimates;
Ouyang and Daganzo~\citep{Ouyang2006} validate the CA scheme against discrete facility
designs, and Ansari et al.~\citep{Ansari2018} survey its application to a broad class
of logistics network-design problems.  Our ruggedness threshold theorem
is a terrain-and-season extension of this CA tradition.

\subsection*{Summary of the gap}

The literature thus supplies: exact and heuristic E-VRP and ELRP algorithms
(on flat networks); physics-based energy models with grade and regen (for
single-vehicle routing); empirical winter range degradation factors; and
robustness frameworks for demand-side uncertainty in facility location.
The gap addressed here is a study that (i) couples arc-level grade \emph{and}
seasonal energy factors into a charger-siting formulation, (ii) measures
the cost of omitting either factor factorially with common random numbers,
and (iii) relates the discrete experiment to a transparent mean-tour
capacity screen.  The present paper addresses that combination while
retaining the limitations of a single-region heuristic study.

%%======================================================================
\section{Problem Setting and Energy Model}
\label{sec:model}

\subsection{Road network and grade data}
\label{sec:network}

The study region is the Trondheim delivery district (Norway), bounded by
63.27--63.52\,N and 10.10--10.90\,E (approximately $45 \times 65$\,km).
The road network is extracted from OpenStreetMap (OSM).  Arc elevation
changes are computed from the Norwegian national LiDAR digital terrain
model (Kartverket NHM; native 1\,m, resampled to 10\,m for the regional
extent), sampled at each arc endpoint in UTM zone 33N (EPSG:25833).  The
resulting network has 296{,}276 directed arcs, of which 291{,}145 have
valid DEM coverage at both endpoints and enter the energy statistics
below.  \Cref{fig:network} shows the study region, the grade structure of
the network, the depot, and the candidate charging sites.

\begin{figure}[ht]
\centering
\includegraphics[width=0.95\textwidth]{network_grade_map.png}
\caption{Trondheim study region.  Road network coloured by absolute
grade (capped at 12\%), with the fleet depot (Sluppen, black triangle),
candidate charging sites (dots), and an illustrative customer draw.
Steep corridors concentrate in the By\aa{}sen hills to the west and the
Gauldalen valley to the south.}
\label{fig:network}
\end{figure}

\emph{Ruggedness} $R$ is the length-weighted mean absolute grade of the
road network (both driving directions, so net elevation cancels):
\begin{equation}
  R = \frac{\sum_a |dz_a|}{\sum_a d_a},
  \label{eq:ruggedness}
\end{equation}
where $dz_a$ and $d_a$ are the signed elevation change and horizontal
length of arc $a$.  For Trondheim, $R = 4.57\%$.  The out-of-sample
transfer cities use Copernicus GLO-30 DSM (30\,m).  Because a surface
model includes buildings and vegetation, its ruggedness values are not
directly comparable with the Trondheim terrain-model estimate; this
measurement limitation is discussed in \cref{sec:predictions}.

\subsection{Physics energy model}
\label{sec:energy}

For an arc of horizontal length $d$\,(m) and signed elevation change
$dz$\,(m), the wheel energy at payload $u$\,(kg above curb mass $m_0$) is:
\begin{equation}
  E_{\text{wheel}}(u) = d\,(k_0 + k_1 u) + \frac{g\,(m_0 + u)\,dz}{J},
  \label{eq:Ewheel}
\end{equation}
where $k_0$ (kWh/m) lumps the rolling and aerodynamic resistance of the
empty vehicle and is anchored to the assumed flat-road consumption of the
hypothetical truck, $k_1$ (kWh/m/kg) is the payload coefficient of rolling
resistance ($k_1 = C_r\,g / J$ with $C_r = 0.0075$),
$g = 9.81$\,m/s$^2$ is gravitational acceleration,
$J = 3.6 \times 10^6$\,J/kWh converts joules to kWh, and the second term
is the gravitational potential energy.  Battery energy consumption is:
\begin{equation}
  E_{\text{batt}}(u) =
  \begin{cases}
    E_{\text{wheel}}(u) / \eta_{\text{drive}} & \text{if } E_{\text{wheel}} \ge 0 \quad \text{(traction)},\\
    E_{\text{wheel}}(u) \cdot \eta_{\text{regen}} & \text{if } E_{\text{wheel}} < 0 \quad \text{(regeneration)},
  \end{cases}
  \label{eq:Ebatt}
\end{equation}
with drive efficiency $\eta_{\text{drive}} = 0.90$ and regeneration
efficiency $\eta_{\text{regen}} = 0.60$ (baseline; sensitivity band
$\eta_{\text{regen}} \in [0.5, 0.7]$ in \cref{sec:sensitivity}).

\emph{Winter scenario.}  At a design temperature of $-3\,^{\circ}$C,
positive battery draw is multiplied by $\alpha_w = 1.35$ and recovered
energy by $\beta_w = 0.80$.  Published studies establish that auxiliary
loads, battery limitations, and reduced recuperation can produce effects of
this order~\citep{Steinstraeter2021,Neubauer2014,Jeffers2016}; they do not
identify these two coefficients for the hypothetical truck used here.
Accordingly, $\alpha_w$ and $\beta_w$ are scenario parameters rather than
an OEM-calibrated winter model.

\subsubsection*{Load-monotone dominance and its failure}
\label{sec:dominance}

Because $E_{\text{batt}}$ involves the grade term $g\,dz/J$, the sign of
$\partial E_{\text{wheel}} / \partial u = d k_1 + g\,dz/J$ can be
negative on downhill arcs steeper than $-k_1 J / g = -C_r \approx
-0.75\%$.  On such arcs a heavier vehicle recovers \emph{more}
charge, invalidating the standard load-monotone dominance rule used in
exact E-VRP-NL label-setting.  Among the 291{,}145 arcs with valid
elevation at both endpoints, 35.9\%
have $\partial E_{\text{batt}} / \partial u < 0$; of all downhill arcs
($dz < 0$), 87.0\% are regen-dominant.  \Cref{app:frvcp} proves that
SoC-function monotonicity and the finite piecewise-linear representation
are nevertheless preserved.  We therefore avoid applying the usual
load-dominance partial order across different delivered-customer subsets.

\subsection{Terrain law and ruggedness threshold}
\label{sec:threshold}

Let $\eff(R)$ denote the fleet-average effective energy consumption
(kWh/km), estimated by simulating the physics model over the real road
geometry of all 2\,km tiles with sufficient road coverage and fitting a
linear terrain law:
\begin{equation}
  \eff(R) = a + b\,R,
  \label{eq:terrainlaw}
\end{equation}
with intercept $a$ pinned to the flat-terrain rate and slope $b$ capturing
the grade-regen interaction.  The fit achieves $R^2 = 0.985$ on
Trondheim 2\,km tiles (\cref{fig:efflaw}).  The intercept is a
physics-derived quantity, not a free parameter; the terrain law is a
spatial summary of the vehicle model applied to real network geometry,
not an empirical vehicle calibration.

\begin{figure}[ht]
\centering
\includegraphics[width=0.78\textwidth]{eff_vs_ruggedness.png}
\caption{Terrain law for the Trondheim region (winter, half payload).
Each point is a 2\,km network tile (marker size proportional to road km);
effective consumption from both-direction traversal of the real road
geometry rises linearly in tile ruggedness $R$.  The pinned fit
$\eff(R) = 0.99 + 10.5\,R$ achieves $R^2 = 0.985$; the dashed line marks
the regional mean $R = 4.6\%$.}
\label{fig:efflaw}
\end{figure}

Winter modifies both intercept and slope: traction increases $a$ and
weakens regeneration (which reduces the fraction of downhill energy
recovered), raising $b$.  For the 14\,t FL-class vehicle:
\begin{equation}
  \eff(R) =
  \begin{cases}
    0.732 + 5.24\,R & \text{summer},\\
    0.989 + 10.47\,R & \text{winter}.
  \end{cases}
  \label{eq:terrainlaw_num}
\end{equation}
The slope roughly doubles in winter ($b_w / b_s \approx 2.0$), reflecting
the asymmetric impact of reduced regen efficiency on rugged terrain.

\subsubsection*{Continuum screening threshold}
\begin{proposition}[Mean-tour ruggedness screen]
\label{thm:threshold}
Let $L$ be the expected tour length (km) under the Daganzo continuum
approximation and $\Qprime = Q(1-r)$ the usable battery capacity with
reserve fraction $r$.  Within this mean-tour approximation, expected tour
energy exceeds usable capacity if and only if
\begin{equation}
  \eff(R) \cdot L > \Qprime
  \quad \Longleftrightarrow \quad
  R > \Rstar = \frac{\Qprime / L - a}{b}.
  \label{eq:Rstar}
\end{equation}
\end{proposition}

\noindent The derivation is given in \cref{app:threshold}.  The threshold has
closed-form comparative statics: $\partial\Rstar/\partial\Qprime = 1/(bL)
> 0$ (larger battery raises the threshold), $\partial\Rstar/\partial L =
-\Qprime/(bL^2) < 0$ (longer tours lower it), and winter lowers $\Rstar$
through both a larger $a$ and a larger $b$.

\subsubsection*{Numerical thresholds (Trondheim)}

\Cref{tab:thresholds} reports $\Rstar$ for each battery $\times$ season
combination at the Trondheim CA tour length $L = 64.7$\,km.  With $R =
4.57\%$, only the 100\,kWh winter case crosses the mean-tour screen.
\Cref{fig:threshold} visualises the threshold as a function of battery
size: the winter curve passes through the regional ruggedness line
between the two candidate pack sizes, which is precisely what makes the
procurement decision an information decision.

\begin{figure}[ht]
\centering
\includegraphics[width=0.78\textwidth]{threshold_phase.png}
\caption{Mean-tour ruggedness screen $\Rstar(Q)$ for summer and winter at the
Trondheim tour length $L = 64.7$\,km.  A region whose ruggedness lies
above the curve has expected tour energy above usable capacity under the
continuum approximation.  The dashed line is
the Trondheim regional $R = 4.6\%$; dotted verticals mark the 100 and
145\,kWh pack options.  Only the 100\,kWh winter combination binds.}
\label{fig:threshold}
\end{figure}

\begin{table}[ht]
\centering
\caption{Ruggedness threshold $\Rstar$ by battery and season.
Trondheim regional ruggedness $R = 4.57\%$.
Binding column indicates $R > \Rstar$.}
\label{tab:thresholds}
\begin{tabular}{llcccc}
\toprule
Battery & Season & $a$ & $b$ & $\Rstar$ & Binds ($R = 4.57\%$) \\
        &        & (kWh/km) & (kWh/km) & & \\
\midrule
145\,kWh & Summer & 0.732 & 5.24 & 24.5\% & No \\
145\,kWh & Winter & 0.989 & 10.47 & 9.8\%  & No \\
100\,kWh & Summer & 0.732 & 5.24 & 12.6\% & No \\
100\,kWh & Winter & 0.989 & 10.47 & \textbf{3.8\%}  & \textbf{Yes} \\
\bottomrule
\end{tabular}
\end{table}

\subsubsection*{Vehicle-class stratification}

The same screen applied to the hypothetical 3.5\,t E-Transit-class van
($m_0 = 2600$\,kg; 89/68\,kWh; $k_0 = 2.2 \times 10^{-4}$; terrain law
$0.26 + 1.4R$ summer / $0.35 + 2.7R$ winter) yields
$\Rstar \in [22\%, 72\%]$ across battery $\times$ season.  None of the
regions evaluated in this study approaches these values, so the screen does
not bind for the modeled van scenarios.  This is a model-based class
comparison, not a universal claim that last-mile vans never require
en-route charging.

\subsection{Vehicle and energy parameters}
\label{sec:params}

\Cref{tab:vehicles} summarises the vehicle configurations used in the
siting experiments.  Both battery variants share the chassis mass and
aerodynamic coefficients.

\begin{table}[ht]
\centering
\caption{Vehicle and energy model parameters.
These are representative hypothetical configurations, not OEM specifications.}
\label{tab:vehicles}
\begin{tabular}{lcc}
\toprule
Parameter & 14\,t FL-class truck & 3.5\,t E-Transit van \\
\midrule
Curb mass $m_0$ (kg)                        & 8,200   & 2,600 \\
Max payload $u_{\max}$ (kg)                 & 5,800   & 1,100 \\
Battery capacity $Q$ (kWh)                  & 145 / 100 & 89 / 68 \\
Vehicle DC power limit (kW)                 & 150     & 180 \\
Rolling+aero coeff.\ $k_0$ (kWh/m)        & $6.0\times10^{-4}$ & $2.2\times10^{-4}$ \\
Payload coeff.\ $k_1$ (kWh/m/kg)          & $2.04\times10^{-8}$ & $2.04\times10^{-8}$ \\
Drive efficiency $\eta_{\text{drive}}$      & 0.90    & 0.90 \\
Regen efficiency $\eta_{\text{regen}}$      & 0.60    & 0.60 \\
Winter traction factor $\alpha_w$           & 1.35    & 1.35 \\
Winter regen factor $\beta_w$              & 0.80    & 0.80 \\
Battery reserve fraction $r$               & 0.10    & 0.10 \\
\bottomrule
\end{tabular}
\end{table}

In the discrete Trondheim experiments, the depot charger is 150\,kW.
Existing candidate sites use the maximum power parsed from OSM when
available; sites with unknown power, including the synthetic coverage-grid
candidates, default to 22\,kW, and all powers are capped by the vehicle
limit.  The 150\,kW power used in the continuum screen is therefore a
strategic reference value, not the power assigned to every discrete site.

%%======================================================================
\section{The Grade+Winter ELRP and Matheuristic}
\label{sec:elrp}

\subsection{Problem formulation}
\label{sec:formulation}

Let $G = (V, A)$ be the road network with node set $V$ and arc set $A$.
The \emph{depot} is node $0 \in V$; the set of customer nodes is
$C \subseteq V$; the set of candidate charging-station locations is
$\mathcal{F} \subseteq V$.  Each customer $i \in C$ has a delivery demand
$q_i$\,(kg), a service time $s_i$\,(min), and a time window $[e_i, l_i]$.
A fleet of identical vehicles with curb mass $m_0$, payload capacity
$u_{\max}$, and battery $Q$ departs from and returns to the depot.

A \emph{site-selection decision} is a subset $F \subseteq \mathcal{F}$,
$|F| \le B$, where $B$ is an upper bound on the number of non-depot sites.
The depot is always
available as a charging site at no cost.  Given $F$, the \emph{routing
problem} seeks a minimum-duration set of routes covering all customers,
where each route must be battery-feasible (SoC remains in $[0, Q]$ at all
times, with reserve $rQ$ maintained) and may visit sites in $F \cup \{0\}$
for en-route charging.

The \emph{\ELRP{}} minimises total route duration (travel time + service
time + charging time + waiting time) over both the site-selection and
routing decisions:
\begin{equation}
  \min_{F \subseteq \mathcal{F},\; |F| \le B} \;
  \min_{\text{routes covering } C \text{ feasible given } F} \;
  \sum_{\text{routes}} \text{duration}(\text{route}, F),
\end{equation}
subject to battery feasibility using the energy model of
\cref{eq:Ewheel,eq:Ebatt}.

The four \emph{planner variants} in the factorial experiment differ only
in which model is used to \emph{solve} the siting problem; all are
\emph{evaluated} against the full (grade+winter) model:
\begin{itemize}
  \item \textbf{Full}: grade+winter model (the reference planner).
  \item \textbf{GOnly}: grade-aware, winter-blind ($\alpha_w = \beta_w = 1$).
  \item \textbf{WOnly}: winter-aware, grade-blind (flat energy model,
    winter factors applied).
  \item \textbf{Blind}: flat model, summer factors only.
\end{itemize}

\subsection{Conservative SoC-grid FRVCP}
\label{sec:frvcp}

Given a fixed ordered sequence of customer visits and candidate charging
stops, the \emph{fixed-route vehicle charging problem} (FRVCP) seeks
the minimum-time charging policy within a restricted charging-detour
neighbourhood.  Following the FRVCP literature
\citep{Montoya2017,Froger2019}, we propagate the earliest clock time at each
point of a uniform SoC grid.  For every customer-to-customer leg, the vehicle
may travel directly or detour through one of the eight open chargers with
the smallest detour time.  Charging amounts are optimized over the grid
using the site's piecewise-linear charging curve.

After each transition, SoC is rounded downward to the grid.  The resulting
route is therefore feasible for the discretized model and conservative with
respect to available charge.  The main experiments use
$\Delta_{\SoC}=0.5$\,kWh, or 0.34--0.50\% of nominal capacity.  This does
not establish a uniform time-domain $\varepsilon$ bound: rounding can alter
a discrete detour or time-window decision.  We therefore describe the method
as a conservative discretization rather than an exact continuous FRVCP
algorithm.  Lexicographic tie-breaking prefers fewer charging stops.

\subsection{ALNS route-pool matheuristic}
\label{sec:alns}

The siting problem is solved by alternating between a route-generation
phase (ALNS heuristic) and a set-partitioning MILP over the accumulated
route pool.

\paragraph{ALNS route generation.}
An adaptive large-neighbourhood search~\citep{Ropke2006} iterates over
destroy operators (random, related, worst-cost, and route removal)
and repair operators (regret-$k$ insertion with $k=2$, greedy insertion).
Candidate insertion positions are prescreened by travel-time detour and the
five best positions are evaluated with the SoC-grid FRVCP.  Operator weights
are adaptive and simulated annealing governs acceptance.  This is route-pool
generation, not column generation: no reduced-cost pricing problem is solved.

\paragraph{Set-partitioning MILP.}
Let $\cR$ be the set of distinct routes accumulated over the ALNS run.
For each route $r \in \cR$, let $x_r \in \{0, 1\}$ be the selection
variable, $F_r \subseteq \mathcal{F}$ the set of charging sites used, and
$z_i$ a binary penalty variable for an uncovered customer.  The
set-partitioning formulation is:
\begin{align}
  \min_{x,\,y,\,z} \quad & \sum_{r \in \cR} c_r\,x_r
    + M\sum_{i\in C}z_i \label{eq:milp_obj}\\
  \text{s.t.} \quad & \sum_{r \in \cR} a_{ir}\,x_r + z_i = 1
    \quad \forall i \in C \label{eq:milp_cover}\\
  & \sum_{s \in \mathcal{F}} y_s \le B \label{eq:milp_budget}\\
  & x_r \le y_s \quad \forall r \in \cR,\; s \in F_r
    \label{eq:milp_site}\\
  & x_r \in \{0, 1\},\; y_s \in \{0, 1\},\; z_i \in \{0,1\}.
    \label{eq:milp_bin}
\end{align}
where $a_{ir} = 1$ if route $r$ serves customer $i$, $c_r$ is the
route duration, $y_s$ indicates whether site $s$ is open, and
$M=10{,}000$ minutes prioritizes coverage over route duration.  The depot
charger is pre-opened ($y_0 = 1$ fixed) and excluded from the budget
count.  Only sites used by selected routes are reported.  The MILP is solved
with Gurobi when available and HiGHS otherwise.  It is exact for the
generated route pool, not for the full location-routing problem.

\subsection{Discretized small-instance reference}
\label{sec:exact}

For benchmark instances with $n \le 10$ customers, a state-space dynamic
program enumerates route customer subsets and propagates states indexed by
(delivered subset, last customer, SoC level), followed by a subset-partition
DP over feasible routes.  Its complexity is of order $3^n$ times the SoC
and transition dimensions.  The reference uses the benchmark-specific
SoC step $Q/200$ and the six smallest-detour charging candidates per leg.
It is optimal conditional on those restrictions.  Because the matheuristic
uses eight charging candidates per leg, the reference is an internal
agreement check rather than a global lower bound or a published BKS.

%%======================================================================
\section{Solver Validation on Standard Benchmarks}
\label{sec:validation}

We exercise the flat special case of the matheuristic on two standard
benchmark sets, 212 instances in total: 92 Schneider (2014) E-VRPTW
instances with $n = 5$--$100$ customers, run through a linear-charging
adapter that maps the PWL charging model to unit-rate linear charging,
and 120 Montoya (2017) E-VRP-NL instances with $n = 10$--$320$
customers, run through a PWL adapter using the published station-type
curves.

On the 44 instances small enough for the discretized reference
($n \le 10$), the matheuristic agrees to numerical tolerance on 42: all
20 Montoya $n = 10$ instances, all 12 Schneider $n = 5$ instances, and
10 of 12 Schneider $n = 10$ instances (\cref{tab:bench}).  The median
gap is 0.00\% and the mean 0.25\%; the only nonzero gaps are rc108C10
(4.48\%) and rc205C10 (6.59\%).  Both are RC-class Schneider instances,
but the current outputs do not isolate whether waiting, route-pool coverage,
or insertion prescreening causes the gaps.

The benchmark file also reports a separate Schneider run with a large
per-route fixed cost to prioritize vehicle count, followed by route duration.
That is not the published vehicle-count-then-distance objective.  The gaps in
\cref{tab:bench} instead compare a pure-duration ALNS run with the
pure-duration discretized DP.  Montoya's total-time objective is closer to
ours.  Since BKS values have not yet been transcribed into the repository,
the remaining 168 runs establish compatibility and scale coverage, not
competitive solution quality.

\begin{table}[ht]
\centering
\caption{Internal benchmark agreement summary (flat special case).
Reference instances have $n \le 10$ and use the discretized state-space
DP; gaps are computed on those instances only.  The two nonzero gaps
are rc108C10 (4.48\%) and rc205C10 (6.59\%).}
\label{tab:bench}
\footnotesize
\setlength{\tabcolsep}{4.5pt}
\begin{tabular}{llrrrrr}
\toprule
Benchmark & $n$ range & \#Inst. & Reference & Numerical & Median gap & Mean gap \\
          &           &         & ($n \le 10$) & agreements & (\%) & (\%) \\
\midrule
Schneider E-VRPTW & 5--100  & 92  & 24 & 22 & 0.00 & 0.46 \\
Montoya E-VRP-NL  & 10--320 & 120 & 20 & 20 & 0.00 & 0.00 \\
\midrule
Total             & 5--320  & 212 & 44 & 42 & 0.00 & 0.25 \\
\bottomrule
\end{tabular}
\end{table}

%%======================================================================
\section{Experimental Design: The Factorial Regret Experiment}
\label{sec:experiment}

\subsection{Instance generation}
\label{sec:instances}

Instances have $N = 100$ customers sampled from a population-weighted
spatial distribution: each candidate customer location is drawn from a
mixture of the SSB 250\,m census grid (weight proportional to population
within 300\,m) and a uniform base component (weight 0.1).  Population is a
transparent residential-delivery proxy, not an observed freight-demand
surface.  Customer demand is sampled uniformly from 150--500\,kg, service
time is 15 minutes, 40\% of customers receive a three-hour time window, and
the route horizon is nine hours.

Ten instances are generated with seeds $s \in \{42, \ldots, 51\}$, each
run under four planner variants and two battery sizes (145\,kWh and
100\,kWh).  The candidate set contains 86 OSM charging locations and
30 terrain-blind 3\,km coverage-grid locations, plus the depot.

\subsection{Common-random-number design}
\label{sec:crn}

Each seed defines one fixed customer instance.  All four planners receive
the same planner seed, and every selected site set is evaluated with the
same three evaluation seeds.  A frozen full-information route-pool snapshot
is copied before each evaluation top-up, preventing evaluation order from
leaking routes between site sets.  Common random numbers reduce avoidable
simulation noise, but planner trajectories still diverge when their energy
models rank moves differently.

\emph{Estimated placement regret} for a misinformed planner $p$ at budget
$B$ and seed $s$ is:
\begin{equation}
  \delta^p_{B,s} = \text{cost}(F^p_{B,s},\,M_{\text{full}})
    - \text{cost}(F^{\text{full}}_{B,s},\,M_{\text{full}}),
  \label{eq:regret}
\end{equation}
where $F^p_{B,s}$ is the site set chosen by planner $p$ and
$F^{\text{full}}_{B,s}$ is the site set returned by the full-information
matheuristic.  Both are evaluated under the full model $M_{\text{full}}$.
The evaluation takes the best of three independent 500-iteration ALNS
top-ups followed by the pool MILP.  Thus the baseline is a matched heuristic
reference, not a proven optimum; individual estimates can be negative
because of residual search error.

The total regret of the \textbf{Blind} planner decomposes as:
\begin{equation}
  \delta^{\text{blind}}_{B,s}
    = \delta^{\text{terrain}}_{B,s}
    + \delta^{\text{season}}_{B,s}
    + \delta^{\text{interaction}}_{B,s},
  \label{eq:regret_decomp}
\end{equation}
where the terrain simple effect is WOnly minus Full, the season simple
effect is GOnly minus Full, and the interaction is the residual.  These are
reference-cell factorial contrasts, not optimization-independent causal
effects.  Their signs may suggest substitution or complementarity, but that
interpretation requires uncertainty intervals that exclude zero.

\subsection{Depot charging counterfactual}
\label{sec:depot}

A free 150\,kW DC charger at the depot is open to all planners and routes
under all scenarios.  The route evaluator may insert a depot detour between
two customer visits, so the experiment measures the incremental value of
non-depot sites conditional on this fallback.

\subsection{Summary statistics and inference}
\label{sec:inference}

For each (battery, budget) cell we report the mean regret and an unadjusted
Student-$t$ 95\% confidence interval across the ten seeds.  For summaries
over $B=2,\ldots,8$, we first average within seed and then form the interval
across ten seed-level means, avoiding pseudoreplication across budgets.  To
assess budget response, we fit a linear slope over $B=2,\ldots,8$ within
each seed and report the mean slope and its $t$ interval.  These intervals
are descriptive and are not multiplicity-adjusted.

%%======================================================================
\section{Results}
\label{sec:results}

\subsection{Battery-dependent information effects}
\label{sec:flip}

\Cref{tab:regret} and \cref{fig:regret_factorial} present mean placement
regret by battery, budget, and component.

\begin{figure}[ht]
\centering
\includegraphics[width=\textwidth]{regret_factorial_n100.png}
\caption{Estimated placement regret (min/day) with unadjusted pointwise
Student-$t$ 95\% confidence bands over ten paired demand seeds.  Component
means change ordering with battery size, but the paired terrain--season
contrasts averaged over $B=2,\ldots,8$ include zero for both packs.
Total regret is larger for 100\,kWh, and neither panel shows evidence of
a decreasing budget trend.}
\label{fig:regret_factorial}
\end{figure}

\paragraph{Total effect.}
Averaged within seed over $B=2,\ldots,8$, blind planning adds
17.5\,min/day for the 145\,kWh pack (95\% CI 0.8--34.2) and
60.8\,min/day for the 100\,kWh pack (16.6--104.9).  Relative to the
full-information baseline, the corresponding means are 0.86\%
(0.04--1.69) and 2.97\% (0.86--5.08).  The total effect is therefore
larger and more operationally relevant for the smaller pack.

\paragraph{Factorial components.}
For 145\,kWh, the seed-averaged terrain and season simple effects are
21.7\,min/day ($-1.1$--44.5) and 8.8\,min/day ($-9.3$--26.9).
For 100\,kWh, they are 15.2\,min/day ($-1.8$--32.1) and
29.5\,min/day (2.3--56.6).  The component means therefore change ordering,
consistent with the threshold screen in \cref{tab:thresholds}.  However,
the paired season-minus-terrain contrasts are $-12.9$\,min/day
($-31.9$--6.1) and 14.3\,min/day ($-17.0$--45.6), respectively.
The present ten-seed experiment supports a shift in relative importance,
not a statistically resolved dominance reversal.  The mean interactions
also change sign, but their seed-level confidence intervals include zero.

\begin{table}[ht]
\centering
\caption{Mean estimated placement regret (min/day, 10 seeds) by battery,
budget, and factorial contrast.  Total = Blind minus Full; Terrain =
WOnly minus Full; Season = GOnly minus Full; Interaction = residual.
The table reports means only; uncertainty is summarized in the text and
shown pointwise in \cref{fig:regret_factorial}.}
\label{tab:regret}
\small
\begin{tabular}{ccrrrr}
\toprule
Battery & $B$ & Total & Terrain & Season & Interaction \\
\midrule
\multirow{8}{*}{145\,kWh}
 & 1 & 14.1 & 14.0 & 11.0 & $-10.9$ \\
 & 2 & 14.4 & 21.4 &  7.1 & $-14.1$ \\
 & 3 & 19.3 & 23.1 & 10.4 & $-14.2$ \\
 & 4 & 15.4 & 19.2 &  6.5 & $-10.3$ \\
 & 5 & 15.4 & 19.2 &  6.5 & $-10.3$ \\
 & 6 & 22.1 & 25.9 & 13.2 & $-16.9$ \\
 & 7 & 17.8 & 21.6 &  8.9 & $-12.7$ \\
 & 8 & 17.8 & 21.6 &  8.9 & $-12.7$ \\
\midrule
\multirow{8}{*}{100\,kWh}
 & 1 & 27.9 &  3.9 & 11.3 & $+12.7$ \\
 & 2 & 56.2 &  6.2 & 27.2 & $+22.8$ \\
 & 3 & 59.6 & 18.9 & 28.4 & $+12.3$ \\
 & 4 & 64.6 & 20.9 & 36.1 & $+7.5$  \\
 & 5 & 58.8 & 15.3 & 24.7 & $+18.8$ \\
 & 6 & 60.0 & 11.4 & 29.9 & $+18.7$ \\
 & 7 & 64.2 & 17.8 & 31.8 & $+14.6$ \\
 & 8 & 62.2 & 15.7 & 28.1 & $+18.4$ \\
\bottomrule
\end{tabular}
\end{table}

\subsection{No decreasing regret trend over the tested budget range}
\label{sec:plateau}

\Cref{tab:regret} shows no decline in total regret from $B=2$ through
$B=8$.  The mean within-seed slope is 0.49\,min/day per additional
budgeted site for 145\,kWh (95\% CI $-0.96$--1.94) and 0.81 for
100\,kWh ($-2.03$--3.64).  Thus the experiment provides no evidence of
repair over this range, but it also does not prove a mathematically flat
response.

The used-site counts identify the immediate mechanism.  At $B=8$, the
Full planner uses on average 2.3 non-depot sites for 145\,kWh and 6.1 for
100\,kWh; the Blind planner uses only 0.2 and 0.9.  Because the constraint
is $|F|\le B$, the blind model is not forced to spend the allowance.  By
omitting grade and winter losses, it sees little incremental value in
en-route charging and leaves most of the budget unused.  The result should
therefore not be interpreted as a test or falsification of facility-location
submodularity.

\subsection{No unserved customers under depot charging}
\label{sec:stranded}

Under every evaluated planner, battery, budget, and seed combination, the
pool MILP leaves zero customers uncovered.  The estimated penalty therefore
appears in travel, charging, and customer-time-window waiting rather than in
the model's unserved-customer penalty.  This does not prove that every
possible operating realization is feasible; it describes the generated
instances and route pools.

\subsection{Prospectively recorded transfer checks}
\label{sec:predictions}

Before fetching Stuttgart and Groningen data, we committed five prediction
groups from the threshold screen to the project repository:

\begin{enumerate}[label=P\arabic*.]
  \item Ruggedness ordering: Trondheim $>$ Stuttgart $>$ Groningen.
    Predicted intervals were 2.5--4.5\% for Stuttgart and 0.1--0.9\%
    for Groningen.
  \item Groningen (flat control): nothing binds for any battery $\times$ season $\times$ class.
  \item Stuttgart boundary: 145\,kWh never binds; 100\,kWh winter within $\pm 1.5$\,pp of $\Rstar$.
  \item The modeled van class remains below its threshold in both cities.
  \item Slope transfer: Stuttgart and Groningen winter slopes $b$ within $\pm 25\%$ of Trondheim's 10.47\,kWh/km.
\end{enumerate}

\Cref{tab:predictions} records the outcomes without collapsing compound
predictions into a single hit rate.  The ordering, Stuttgart interval,
binding classifications, van classification, and slope-transfer checks
were consistent with the recorded expectations.  Groningen ruggedness was
2.2\%, outside the predicted 0.1--0.9\% interval, and its closest
$\Rstar/R$ ratio was 4.1 rather than the predicted minimum of 5.
Copernicus GLO-30 is a DSM whereas Trondheim uses a LiDAR terrain model, so
surface artefacts are a plausible explanation; without a matched Groningen
DTM comparison, they are not a demonstrated causal diagnosis.

\begin{table}[ht]
\centering
\caption{Prospectively recorded transfer-check outcomes for Stuttgart
and Groningen.  Predictions were committed to the repository before the
city-data fetch; the timestamped record is included in the supplement.}
\label{tab:predictions}
\small
\begin{tabular}{>{\raggedright\arraybackslash}p{2.4cm}
                >{\raggedright\arraybackslash}p{5.1cm}
                >{\raggedright\arraybackslash}p{4.7cm}}
\toprule
Pred.\ & Statement & Outcome \\
\midrule
P1 (ordering) & $R_{\text{Trd}} > R_{\text{Stgt}} > R_{\text{Gron}}$ & Consistent: 4.6\% > 4.2\% > 2.2\% \\
P1 (Stuttgart) & $R \in [2.5\%, 4.5\%]$ & Consistent: 4.2\% \\
P1 (Groningen) & $R \in [0.1\%, 0.9\%]$ & Outside interval: 2.2\% \\
P2 & Groningen: no binding; closest $\Rstar/R\ge5$ & No binding, but ratio 4.1 \\
P3 & Stuttgart 100\,kWh winter $|\Rstar - R| \le 1.5$\,pp & Consistent: 1.2\,pp (binding side) \\
P4 & Modeled van does not bind & Consistent: minimum van $\Rstar = 19.7\%$ \\
P5 & $b$ within $\pm 25\%$ of 10.47 & Consistent: Stgt 10.5, Gron 8.5 \\
\bottomrule
\end{tabular}
\end{table}

\subsection{Sensitivity analysis}
\label{sec:sensitivity}

\paragraph{Regeneration efficiency band.}
\Cref{tab:sensitivity} shows the threshold $\Rstar$ for each
$\eta_{\text{regen}} \in \{0.5, 0.6, 0.7\}$.  For the small100 winter
case, $\Rstar$ ranges from 3.6\% to 4.2\% across the band, remaining below
the regional $R = 4.57\%$ throughout.  The binding verdict is invariant
across the entire sensitivity band.  This sensitivity analysis concerns the
mean-tour threshold only; it does not establish robustness of the discrete
factorial components.

\begin{table}[ht]
\centering
\caption{Threshold sensitivity to regeneration efficiency
$\eta_{\text{regen}}$.  All three values bind for small100 winter
($\Rstar < R = 4.57\%$); no other scenario binds.}
\label{tab:sensitivity}
\begin{tabular}{llccc}
\toprule
Battery & Season & $\eta_{\text{regen}} = 0.5$ & $0.6$ & $0.7$ \\
\midrule
145\,kWh & Summer & 20.5\% & 24.5\% & 30.5\% \\
145\,kWh & Winter &  9.1\% &  9.8\% & 12.1\% \\
100\,kWh & Summer & 10.5\% & 12.6\% & 15.6\% \\
100\,kWh & Winter & \textbf{3.6\%} & \textbf{3.8\%} & \textbf{4.2\%} \\
\bottomrule
\end{tabular}
\end{table}

%%======================================================================
\section{Discussion and Conclusions}
\label{sec:discussion}

\subsection{Procurement-vs-information tradeoff}
\label{sec:procurement}

The strongest procurement result concerns the total information effect,
not the ranking of its components.  Relative blind-planning regret is
approximately 0.9\% for 145\,kWh and 3.0\% for 100\,kWh over the tested
budgets.  The mean-tour screen offers a coherent explanation: 100\,kWh
in winter crosses the aggregate capacity boundary, whereas 145\,kWh
does not.  A smaller pack can therefore increase the value of collecting
and representing environmental information in the siting study.

This is not a complete procurement calculation.  Battery capital cost,
payload displacement, degradation, energy price, and charging-infrastructure
cost are outside the objective.  The results can inform such a calculation
but do not by themselves establish that either pack is economically
preferred~\citep{Basma2022}.

\subsection{Why additional allowance is unused}
\label{sec:plateau_why}

The blind planner's low used-site count is consistent with a demand-
suppression mechanism.  On flat summer roads, the depot fallback makes most
non-depot charging routes unattractive, so increasing an upper bound does
not force the planner to install sites that its own model does not value.
The full model perceives more charging demand and uses more of the available
allowance.  Distinguishing this mechanism from site-rank misspecification
would require a separate exactly-$B$ experiment or a construction-cost
objective; neither is included here.

\subsection{Limitations}
\label{sec:limitations}

\emph{Single-region depth.}  The siting experiment is conducted only in
Trondheim; Stuttgart and Groningen serve as threshold transfer checks, not
full siting experiments.  Results may not generalise quantitatively to
regions with very different grade distributions.

\emph{Hypothetical vehicle classes.}  The 14\,t FL-class and 3.5\,t van
parameters define plausible scenarios but are not calibrated against OEM
test data or truck telemetry.  The energy model is used for relative
comparison rather than absolute prediction.

\emph{Design-temperature winter model.}  A single $-3\,^{\circ}$C design
temperature and two assumed multipliers are used.  They are informed by
passenger-EV evidence but are not estimated for the modeled truck.

\emph{Demand proxy and charging supply.}  Population is not observed freight
demand.  OSM charger power is incomplete, synthetic candidates default to
22\,kW, and charging queues, station capacities, construction costs, grid
constraints, and charger availability are not modeled.

\emph{Heuristic and discretization error.}  The full-information baseline is
not a proven optimum.  The route evaluator uses a downward-rounded SoC grid
and a top-$k$ charging neighbourhood.  The internal reference differs in
$k$, and the benchmark file does not yet include published BKS values or
runtime results.  The two nonzero small-instance gaps are 4.48\% and 6.59\%.

\emph{Transfer elevation data.}  Trondheim uses a LiDAR terrain model,
whereas Stuttgart and Groningen use a 30\,m surface model.  A matched
DTM/DSM comparison is needed before attributing the Groningen miss to a
specific bias magnitude.

\emph{Statistical precision.}  Ten paired seeds resolve the total blind
effect but not the terrain-versus-season contrast or interaction.  The
pointwise intervals are exploratory and are not multiplicity-adjusted.

\subsection{Forward outlook}
\label{sec:forward}

The threshold relation makes $\Rstar(Q, L, \text{season})$ computable from
topographic and road data once vehicle and climate scenarios are supplied.
A natural follow-up study maps
$\Rstar$ across a broader portfolio of cities to classify regions into:
mean-tour binding and non-binding cases, then tests those classifications
with observed freight demand and discrete siting replications.  A stronger
strategic model would endogenize facility density, construction cost, and
queueing rather than treating the threshold as an optimization result.

\subsection{Conclusions}
\label{sec:conclusions}

In this Trondheim experiment, omitting both terrain and winter information
increases estimated fleet route duration by 17.5\,min/day with a 145\,kWh
pack and 60.8\,min/day with a 100\,kWh pack over budgets $B=2,\ldots,8$.
The component means shift from a larger terrain effect at 145\,kWh to a
larger winter effect at 100\,kWh, but the paired contrasts remain
statistically unresolved.  The defensible finding is therefore
battery-dependent information value, not a universal dominance flip.

Estimated regret does not decline over the tested budget range because the
blind planner leaves most of the site allowance unused.  With depot charging,
no customer is left uncovered, so the measured consequence is additional
route duration rather than modeled service failure.  The continuum threshold
helps interpret why the small-pack winter case is sensitive, while its
transfer checks also show the need for comparable elevation instruments.
For practice, battery sizing and model fidelity should be evaluated jointly;
the present results quantify one operational part of that decision.

%%======================================================================
\section*{Acknowledgements}
[To be completed.]

\section*{Data availability}
The code, processed result tables, and instructions for rebuilding the
open-data inputs will be archived in a versioned repository.  Repository URL,
software licence, and data licences must be inserted before submission.

\section*{CRediT author statement}
[To be completed.]

\section*{Funding}
[To be completed.]

\section*{Declaration of competing interest}
The author declares no known competing financial interests or personal
relationships that could have appeared to influence the work reported here.

%%======================================================================
%% BIBLIOGRAPHY
%%======================================================================
\bibliographystyle{elsarticle-harv}
\bibliography{references}

%%======================================================================
%% APPENDICES
%%======================================================================
\appendix

\section{SoC-Function Monotonicity under Regeneration}
\label{app:frvcp}

This appendix proves that the piecewise-linear SoC-function representation
is preserved under the grade+regen energy model of \cref{sec:energy}.

\begin{lemma}[SoC-function monotonicity under regen]
\label{lem:monotone}
Let $f: [0, Q] \to \mathbb{R} \cup \{+\infty\}$ be a non-increasing PWL
function representing the minimum continuation time from a node $w$.
Prepending arc $(u, w)$ with battery energy $E_{uw}$
(which may be negative for downhill arcs) produces the composition:
\begin{equation}
f'(y)=
\begin{cases}
f\bigl(\min\{Q,\,y-E_{uw}\}\bigr)+t_{uw},
  & \min\{Q,\,y-E_{uw}\}\ge rQ,\\
+\infty, & \text{otherwise},
\end{cases}
\end{equation}
where $y$ is the arrival SoC at $u$ and $t_{uw}$ is the arc travel time.
Then $f'$ is non-increasing and PWL.
\end{lemma}

\begin{proof}
Let $h(y) = \min(Q, y - E_{uw})$.  This function is non-decreasing and
piecewise-linear in $y$ (with a single breakpoint at $y = Q + E_{uw}$ where
the ceiling binds for downhill arcs with $E_{uw} < 0$).  The composition
$f \circ h$ is non-increasing because $f$ is non-increasing and $h$ is
non-decreasing.  Adding the constant $t_{uw}$ preserves both properties.
Restricting the finite-valued portion to $h(y)\ge rQ$ adds at most one
feasibility breakpoint and preserves extended-real monotonicity.
Together with the ceiling breakpoint of $h$, the number of breakpoints
therefore increases by at most two, so the PWL representation remains finite.
\end{proof}

\begin{remark}
For downhill arcs ($E_{uw} < 0$), the ceiling $\min(Q, \cdot)$ introduces a
\emph{saturation piece}: when the arrival SoC satisfies
$y \ge Q + E_{uw}$,
the battery reaches the ceiling and the excess regen is wasted.  The
minimum continuation time is therefore $f(Q) + t_{uw}$ for all such $y$,
creating a flat segment in $f'$.  This is the only structural difference
from the flat-network case.
\end{remark}

\section{Mean-Tour Threshold Derivation}
\label{app:threshold}

\begin{proof}[Derivation of \cref{thm:threshold}]
Under the Daganzo continuum approximation, the expected energy consumed on
a tour of length $L$\,(km) with fleet-average effective consumption
$\eff(R) = a + bR$\,(kWh/km) is $\eff(R) \cdot L$.  A tour is
below the mean-tour capacity screen if and only if this expected energy does
not exceed the usable capacity $\Qprime = Q(1-r)$:
\begin{equation*}
  \eff(R) \cdot L \le \Qprime
  \quad \Longleftrightarrow \quad
  (a + bR) L \le \Qprime
  \quad \Longleftrightarrow \quad
  R \le \frac{\Qprime/L - a}{b} = \Rstar.
\end{equation*}
The mean-tour screen is crossed iff $R > \Rstar$.  This algebra does not
claim that every discrete route binds on the same side of the threshold.
The comparative statics follow by differentiation.  Note that
$\partial\Rstar/\partial Q'
= 1/(bL) > 0$ because $b, L > 0$; and $\partial\Rstar/\partial b =
-(Q'/L - a)/b^2 < 0$ iff $a < Q'/L$, which holds when the flat-terrain
consumption is less than the battery-per-km ratio---a necessary condition
for any feasible tour.
\end{proof}

\section{Conservative SoC Discretization}
\label{app:eps_exact}

The SoC-grid DP discretises the SoC axis into steps of $\Delta_\SoC =
0.5$\,kWh in the main experiments.  After each transition, reachable SoC
is rounded down to the next grid point.  Consequently, the stored state
never credits energy that the continuous transition did not provide, and a
route returned as feasible has a feasible counterpart under the same
discrete detour sequence with at least the represented SoC.

This argument establishes conservatism in represented charge, not a global
time-error bound.  Up to one grid step can be lost at every transition, and
near a time-window or detour-choice boundary that loss can change the
discrete decision.  A formal continuous-optimum bound would require
additional regularity assumptions and a grid-convergence study.  The
benchmark comparison therefore uses the term \emph{discretized reference}
and matches the SoC resolution between the reference and the reported
pure-duration run.

\end{document}
